% Автоматически перенесено из ВВЕДЕНИЕ_ВЕРСИЯ_3.docx

История средневекового права долгое время изучалась преимущественно через институциональную историю, развитие нормативных систем и становление правовых памятников. В центре внимания исследователей находились кодексы, сборники обычаев, городские статуты и иные формы письменной фиксации нормы. Применительно к Каталонии XII--XIII вв. особое значение традиционно придавалось «Барселонским обычаям» и связанным с ними памятникам местного права, которые рассматривались прежде всего как свидетельства формирования региональной правовой традиции, взаимодействия римского, канонического, вестготского и локального права\footnote{\emph{Берман }\emph{Г.Дж}\emph{. }Западная традиция права: эпоха формирования / пер. с англ. Н.Р. Никоновой, Н. Н. Деева. 2-е изд. М.: Изд-во МГУ; Издательская группа ИНФРА-М---НОРМА, 1998. С. 26.}.

Такой подход позволял реконструировать эволюцию правовых институтов, выявить политические и социальные условия развития каталонского обычного права, проследить соотношение местного обычая с более широкими правовыми традициями латинского Запада.

Однако институционально-нормативная перспектива неизбежно смещала исследовательский акцент к тому~уровню правовой культуры\footnote{В историографии понятие правовой культуры первоначально оформилось в рамках социологии права и связывалось с ценностными ориентациями общества по отношению к правовой системе: \emph{Friedman}\emph{~L.M.} Legal Culture and Social Development // Law \& Society Review. 1969. Vol. 4, № 1. P. 29--44; \emph{Idem.} The Legal System: A Social Science Perspective. New York: Russell Sage Foundation, 1975. В сравнительно-правовых исследованиях близкая проблематика развивалась через понятие правовой традиции, позволявшее рассматривать право не только как систему норм, но и как совокупность устойчивых представлений о природе права, его роли, способах применения, изучения и передачи: \emph{Merryman}\emph{~}\emph{J.H.}~The Civil Law Tradition: An Introduction to the Legal Systems of Western Europe and Latin America. Stanford: Stanford University Press, 1969. Для анализа сосуществования различных нормативных порядков и взаимодействия официального права, обычая и социальных практик принципиальное значение получили исследования правового плюрализма: \emph{Merry}\emph{~S.} E. Legal Pluralism // Law \& Society Review. 1988. Vol. 22, № 5. P. 869--896. В более поздней литературе понятие правовой культуры стало использоваться осторожнее: исследователи подчёркивали необходимость учитывать его связь с социальными практиками, правовым воображением, профессиональным знанием и формами правовой коммуникации: \emph{Nelken}\emph{ D.} Using the Concept of Legal Culture // Australian Journal of Legal Philosophy. 2004. Vol. 29. P. 1--26; \emph{Cotterrell R.} Law, Culture and Society: Legal Ideas in the Mirror of Social Theory. Aldershot: Ashgate, 2006. В медиевистике близкая исследовательская перспектива позволила рассматривать право как часть интеллектуальной, текстовой и социальной жизни средневекового общества, где обычай, авторитет, память, письменная форма и судебная практика не отделены друг от друга: \emph{Берман }\emph{Г.Дж}\emph{. }Западная традиция права: эпоха формирования / пер. с англ. Н.Р. Никоновой, Н. Н. Деева. 2-е изд. М.: Изд-во МГУ; Издательская группа ИНФРА-М---НОРМА, 1998; \emph{Grossi~P}\emph{.} L’ordine giuridico medievale. Roma; Bari: Laterza, 1995; \emph{Wormald~P. }Legal Culture in the Early Medieval West: Law as Text, Image and Experience. London: Hambledon Press, 1999.}, где право представало как система правил, закреплённых в письменной форме и санкционированных властью\footnote{См. об историко-антропологическом и правовом подходе: \emph{Блок~М.} Апология истории, или Ремесло историка. М., 1986. С. 17--31; \emph{Варьяш}\emph{~}\emph{О.И. }История права и антропология права // \emph{Варьяш}\emph{~}\emph{О.И.} Пиренейские тетради: право, общество, власть и человек в средние века. М.: Наука, 2006. С. 28--29; \emph{Репина}\emph{~}\emph{Л.П., Зверева}\emph{~}\emph{В.В., Парамонова~М.Ю. }История исторического знания. М.: Дрофа, 2004. С. 245.}.

Правовая культура в настоящем исследовании понимается как исторически конкретная среда регулирования общественных отношений, в которой взаимодействовали нормы, представления о праве, способы правового действия и практики их применения. Применительно к Каталонии XII--XIII вв. исследовать эту среду позволяют прежде всего записи обычного права и актовый материал, которые дают возможность проследить передачу норм и формул от одного источника к другому, способы выражения правового содержания на латинском и старокаталанском языках и организацию судебной процедуры.\footnote{\emph{Варьяш}\emph{ О.И. }Язык средневекового права // \emph{Варьяш}\emph{ О.И. }Пиренейские тетради: право, общество, власть и человек в средние века. М.: Наука, 2006. С. 90.}.

В соответствии с этим правовая культура рассматривается в диссертации через три взаимосвязанных аспекта. Первый связан с составом правовых текстов и текстуальными связями между латинской правовой традицией, каталонскими сводами обычного права и актовым материалом. Второй охватывает способы выражения нормы и правового действия: речевые формулы, прямую речь, взаимодействие латыни и старокаталанского языка, заимствования, цитирование и перевод. Третий составляет организация судебного доказательства --- соотношение ордалии, судебного поединка, клятвы, свидетельских показаний и других видов доказательства.

\textbf{Научная проблема} настоящей диссертации состоит в том, чтобы рассмотреть каталонские своды обычного права XII--XIII вв. и актовый материал в совокупности как тексты, в которых отражены способы бытования права в социальной среде. В центре исследования находится взаимодействие нормы, языка, процедуры и человека в различных ситуациях.

\textbf{Актуальность исследования} определяется необходимостью более полного изучения средневекового права как культурной и антропологической системы. В современной медиевистике всё большее значение приобретают вопросы, связанные с историей письменности, устной коммуникации, судебного ритуала, правовой антропологии и языковых форм социального порядка\footnote{См. например, работу нашей коллеги с добротным историографическим обзором: \emph{Сафронова}\emph{~}\emph{Д.П.} Устная речь в правовой и документальной практике Леонского королевства X--XI веков // Диалог со временем. 2020. № 72. С. 125.}.

Каталонский материал особенно показателен для такого исследования, поскольку позволяет сопоставить одни и те же нормы, термины и формулы в различных типах источников, проследить их передачу и переработку в латинской и старокаталанской языковой среде и установить их место в конкретной правовой процедуре. Вычислительное сопоставление выявляет разные степени текстуальной близости между памятниками. Анализ переводов показывает устойчивые соответствия и смысловые трансформации, в то время как, процессуальные нормы позволяют определить, как различались функции клятвы, свидетельства, репутации, ордалии, судебного поединка и документа. Такое соединение текстологического, языкового и классического исторического анализа позволяет перейти от изучения отдельных норм к исследованию связей между способами их передачи, выражения и применения.

Актуальность диссертации связана также с недостаточной изученностью каталонских сводов обычного права в отечественной медиевистике. Отдельные сюжеты, связанные с историей Каталонии, «Барселонскими обычаями», «Обычаями Тортосы» и правовой культурой Пиренейского полуострова, затрагивались в более широком контексте истории западноевропейского права, городской истории и истории Средневековья.

\section*{Степень разработанности темы}
\addcontentsline{toc}{section}{Степень разработанности темы}

Проблематика диссертации находится на пересечении нескольких исследовательских направлений: истории средневекового права, источниковедения и текстологии правовых памятников, социальной истории Каталонии, исторической антропологии, социолингвистики и истории судебной процедуры.

Наиболее разработанным направлением остаётся изучение каталонских сводов обычного права с позиций классического источниковедения. Уже в исследованиях конца XIX --- первой половины XX в. были поставлены основные вопросы о происхождении, датировке, составе и рукописной традиции «Барселонских обычаев». В дальнейшем изучение их редакций, глосс, комментариев и компилятивной традиции получило развитие прежде всего в работах А. Иглесии Феррейроса. Отдельная исследовательская традиция сложилась вокруг «Обычаев Льейды», «Обычаев Тортосы», «Кутюмов Перпиньяна» и связанных с ними локальных памятников.

Другой круг исследований связан с социальной и политической историей Каталонии. Работы П. Боннасси, Т. Н. Биссона, П. Фридмана, Л. То Фигераса, Ф. Сабате-и-Куруя и других исследователей позволяют рассматривать обычное право в контексте конкретных отношений власти, собственности и зависимости, определявших положение участников правового взаимодействия, границы сеньориальных и городских юрисдикций и условия возникновения и разрешения конфликтов.

Значительный вклад в изучение документальной культуры и языковой ситуации в Каталонии обозначенного периода внес М.~Зиммерман. В более широком европейском контексте важны исследования соотношения устного и письменного, грамотности, вернакуляризации и роли документа в средневековом обществе. С этим направлением связаны исследования мультилингвизма и диглоссии, латинской и романской правовой лексики, а также перевода правовых текстов с латыни на вернакуляры.

Для анализа судебного доказательства, клятвы, ордалии, судебного поединка, «доброй славы», свидетельства и репутационных способов доказательства важны исследования, посвящённые судебным ритуалам и трансформации процедуры в XII--XIII вв. Работы Р.~Бартлетта, Д.~Бартелеми, Ж.~Тери и других медиевистов позволяют рассматривать клятву, ордалию, судебный поединок и писанный документ как самостоятельные способы доказательства.

В отечественной традиции значимое место занимают работы А.Я.~Гуревича и О.И.~Варьяш, в которых средневековые правовые источники рассматриваются в связи с представлениями о человеке, статусе, ответственности и социальном порядке. В более широком контексте значение имеют труды М.~Блока, Г.Дж. Бермана, Э. Коэн, М. Фуко и других авторов, позволившие рассматривать право наряду с нормами и институтами через социальные практики, ритуалы и формы разрешения конфликтов.

В отечественной историографии правовая культура Пиренейского полуострова изучалась неравномерно. Дореволюционные и советские исследования затрагивали прежде всего социально-экономическую историю, крепостное право, феодальные отношения и отдельные аспекты истории Испании. В современной отечественной традиции важное значение имеют работы Г.А.~Поповой, И.И.~Варьяш и других исследователей, посвящённые правовой культуре, вестготскому и кастильскому правовому наследию, языку средневекового права, городским обществам и переводам правовых текстов.

Таким образом, достаточно подробно изучены история формирования и рукописная традиция отдельных каталонских сводов обычного права, социально-политический контекст их возникновения, документальная культура Каталонии, языковая ситуация и отдельные формы судебного доказательства. В специальном исследовании до сих пор не сопоставлялись структура текстуальных связей между сводами обычного права и актовым материалом, особенности латинского и старокаталанского словоупотребления и перевода и организация судебного доказательства. Именно соединение этих трёх направлений позволяет поставить вопрос о взаимосвязях между составом и переработкой правовых памятников, особенностями их языка и зафиксированными в них процедурами.

\textbf{Источниковая база исследования }включает три основные группы: своды каталонского обычного права и связанные с ними нормативные акты; актовый материал, фиксировавший конкретные правовые действия и конфликты; тексты римской, канонической, вестготской и книжной традиции, привлекаемые для сопоставления. Различия между этими источниками по происхождению, жанру, языку и функции позволяют сопоставлять формулирование нормы в своде обычного права, её текстовые параллели в других памятниках и способы оформления конкретного правового действия.

Центральное место в исследовании занимают «Барселонские обычаи» и «Обычаи Тортосы». «Барселонские обычаи» используются как основной источник для выявления текстуальных связей между каталонскими сводами обычного права, более ранними латинскими памятниками и актовым материалом. Сопоставление латинского текста со старокаталанскими версиями позволяет отдельно исследовать терминологические соответствия, переводческие решения и изменения, возникавшие при передаче текста на другом языке.

«Барселонские обычаи» рассматриваются как один из основных памятников раннего периода становлений каталонской правовой культуры. Латинский текст позволяет проследить формирование нормативного языка и связь каталонского материала с более широкими правовыми традициями латинского Запада. Старокаталанская версия памятника привлекается для анализа переводческих решений, терминологических соответствий и смысловых сдвигов.

«Обычаи Тортосы» особенно значимы для исследования благодаря объёму и внутреннему разнообразию материала. В работе учитываются две редакции памятника --- 1272 г. и 1276--1279 гг. В них представлены прямая речь и формализованный судебный диалог, речевые формулы, дефиниции, латинские включения в старокаталанский текст и развитая правовая терминология. Одновременно памятник содержит подробную регламентацию судебной клятвы, репутационных критериев, свидетельских показаний и письменного доказательства.

Для сопоставления с этими памятниками привлекаются другие сборники городского и местного обычного права,: «Обычаи Льейды» 1228 г., «Кутюмы Перпиньяна» 1242 г., «Обычаи Тарреги» 1242 г., «Обычаи Орты» 1296 г., «Обычаи Валь-д’Арана», «Обычаи Миравета» 1319 г., королевская привилегия «Знатные мужи Барселоны признали…» (Recognoverunt proceres), пожалованная Барселоне в 1284 г., а также королевские нормативные акты -- две прагматики Жауме II 1295 и 1301 гг. Сопоставление этих памятников позволяет определить характер текстуальной близости между отдельными сводами, выявить устойчивые и локально варьирующиеся формулы и термины и проследить переработку норм применительно к различным городам и территориям.

Особую группу образуют «Обычаи Тарреги», «Обычаи Орты» и «Обычаи Миравета», текстуально связанные с «Обычаями Льейды». Их сопоставление позволяет проследить сохранение и переработку отдельных норм и формул в памятниках, созданных для разных локальных условий, в том числе на территориях Новой Каталонии и во владениях духовно-рыцарских орденов. Совокупность этих памятников демонстрирует, каким образом одни и те же нормы воспринималась, перерабатывалась и приспосабливалась к различным локальным условиям, в том числе на территориях Новой Каталонии и во владениях духовно-рыцарских орденов. «Обычаи Миравета» выходят за основные хронологические рамки исследования и относятся уже к началу XIV в., однако их привлечение необходимо для анализа рецепции норм «Обычаев Льейды» в~XIII~в.

«Кутюмы Перпиньяна» важны как памятник городского права Северной Каталонии и Руссильона, позволяющий проследить соотношение местного обычая, римского права и правовых традиций Южной Франции.

Привилегия «Знатные мужи Барселоны признали…» (Recognoverunt proceres) рассматривается как памятник барселонского городского права конца XIII в. Наличие латинского текста и ранней старокаталанской версии позволяет привлекать её к исследованию правовой терминологии и перевода; содержащиеся в ней формулы имущественных распоряжений, нормы о завещании и судебной процедуре --- к анализу речевых формул, письменного удостоверения и изменений процессуальной регламентации. «Обычаи Валь-д’Арана» и прагматики Жауме II привлекаются как дополнительные сравнительные материалы при анализе позднейшего закрепления отдельных норм и процедур.

Важную часть источниковой базы составляет актовый материал: протоколы судебных заседаний, судебные решения, свидетельские показания, жалобы, соглашения о мирном урегулировании спора, акты признания вины или отказа от прав, а также документы, связанные с обеспечением и исполнением судебных решений (обязательства явки в суд, поручительства, продажа заложенного имущества и осмотр спорных границ). В работе используются комплекс документов IX--XI вв. и XII в, опубликованный в серии «Каталонские юридические тексты» (2018 и 2024 гг.). Актовый материал используется как самостоятельный источник по истории правовой культуры. Грамоты дают возможность выявить устойчивые клише, встречающиеся также в нормативных памятниках, проследить употребление правовой лексики и установить, каким образом в документах фиксировались свидетельские показания, клятва, признание права, отказ от притязаний и другие элементы разрешения спора.

Для выявления возможных источников текстуальных параллелей, цитат и понятийных соответствий привлекаются также памятники римской, канонической, вестготской и книжной традиции: тексты Священного Писания (прежде всего Евангелий), Corpus Iuris Civilis, V книга «Этимологий» Исидора Севильского «О законах и временах»; «Вестготская правда» («Книга приговоров»), «Составленные Петром извлечения из римских законов», а также отдельные памятники канонического права и связанные с ними компилятивные тексты. Эти источники важны как правовая, книжная и дидактическая среда, в которой формировались и функционировали основные памятники каталонского обычного права.

Особое значение среди них имеет «Вестготская правда», поскольку она сохраняла значение для каталонской правовой культуры, использовалась в судебной и образовательной практике, цитировалась и переосмысливалась в «Барселонских обычаях». Corpus Iuris Civilis, канонические тексты и «Составленные Петром извлечения из римских законов»\footnote{Составленные Петром извлечения из римских законов / пер., ввод. ст. и коммент. Д.Ю.~Полдникова; под ред. Л.Л.~Кофанова // Древнее право. Ius antiquum. 2003. № 1 (11). С.~216--262; 2003. № 2 (12). С. 212--231; 2004. № 1 (13). С. 188--207.} позволяют выявлять следы учёного права, прямые и косвенные заимствования, устойчивые формулы и понятийные переклички. «Этимологии» Исидора Севильского и евангельские тексты используются для выявления более широкого книжного и правового словаря, знакомого средневековым легистам.

Совместное использование этих групп источников позволяет сопоставить текстуальные связи между памятниками, особенности их латинских и старокаталанских версий и перевода и зафиксированные в них способы судебного доказательства. Именно такое сопоставление делает возможным исследование тех взаимосвязей внутри каталонской правовой культуры, которые составляют предмет настоящей диссертации.

\textbf{Объектом исследования} является правовая культура Каталонии XII--XIII вв. \textbf{Предмет исследования }составляют текстуальные взаимодействия между каталонскими памятниками обычного права, формы языковой фиксации правовых действий и перевода правовой терминологии, а также система судебного доказательства, рассматриваемые в их взаимосвязи.

\textbf{Цель исследования} заключается в характеристике правовой культуры Каталонии XII--XIII вв. через выявление взаимосвязей между текстуальной передачей обычного права, способами выражения нормы и организацией судебного доказательства. Для достижения поставленной цели необходимо решить следующие\textbf{ }\textbf{исследовательские }\textbf{задачи:}

\begin{enumerate}
\setlength{\itemsep}{0pt}
\item выявить и исторически интерпретировать текстуальные связи между источниками латинской правовой традиции, каталонскими сводами обычного права и актовым материалом с применением количественных методов анализа текста.
\item охарактеризовать способы письменного представления словесного правового действия в каталонских источниках: прямую речь, речевые формулы и глаголы речи;
\item исследовать взаимодействие латинского и старокаталанского языков в правовых источниках, включая иноязычные элементы, заимствования, цитирование и способы введения и объяснения правовой терминологии;
\item сопоставить латинский текст «Барселонских обычаев» с его старокаталанскими версиями и выявить устойчивые терминологические соответствия, вариативность перевода и основные способы лексической, грамматической и синтаксической переработки нормы;
\item определить место и функции ордалии, судебного поединка, судебной клятвы, репутации, свидетельских показаний и письменного документа в системе судебного доказательства и проследить изменение их соотношения в XII--XIII вв.
\end{enumerate}

\textbf{Х}\textbf{ронологические рамки исследования} охватывают XII--XIII вв. Именно к этому времени относятся составление основного латинского текста «Барселонских обычаев» и запись ряда городских и локальных сводов обычного права, составляющих основную источниковую базу исследования. В XIII в. появляются «Обычаи Льейды», «Кутюмы Перпиньяна», «Обычаи Тарреги», «Обычаи Тортосы», «Обычаи Орты» и королевская привилегия «Знатные мужи Барселоны признали…» (Recognoverunt proceres).

При этом датировки отдельных источников, использованных нами, выходят за границы хронологических рамок: для сравнительного анализа в отдельных случаях привлекаются грамоты начиная с IX в. Для понимания более раннего правового наследия, прежде всего, это «Вестготская правда». Кроме того, для изучения позднейшей передачи и переработки установлений обычного права нами были привлечены отдельные памятники XIV в., в частности «Обычаи Валь-д’Арана» 1311 г. и «Обычаи Миравета» 1319 г.

\textbf{Географические рамки исследования }ограничены территорией средневековой Каталонии --- областью, включавшей графства и земли, связанные с политическим и правовым пространством\footnote{Для исследования особенно значимо понятие правового пространства, позволяющее рассматривать территорию не только как географическую область, но и как совокупность пересекающихся юрисдикций, норм, статусов и механизмов правового взаимодействия: \emph{Варьяш}\emph{ И.И.} Сарацины под властью арагонских королей: исследование правового пространства XIV в. СПб.: Евразия Clio, 2016. С. 7--8.} графства Барселонского и, шире, Арагонской Короны, объединившей королевство Арагон, Каталонский принципат, позднее --- королевства Майорку и Валенсию\footnote{Подробнее см. раздел «Каталонские графства в XI--XIII вв.».}. Каталония, с одной стороны, была связана с провансальскими и североитальянскими землями, с другой --- находилась в пространстве Пиренейского полуострова и взаимодействовала с исламским миром. Это проявлялась в постоянном взаимодействии местного обычая, права общин иноверцев, римско-канонических норм и наследия вестготской традиции в рамках локальной правовой культуры. Кроме того, пограничное положение Каталонии повлияло и на языковые особенности правовых памятников, составленных на территории региона.

Специфика исследования требует сочетания методов, позволяющих работать с правовыми источниками одновременно как с историческими памятниками, текстами на латинском и старокаталанском языках и свидетельствами судебной практики. \textbf{Методологическую основу }\textbf{исследования} составляют историко-правовой, источниковедческий, историко-лингвистический и историко-антропологический подходы; для выявления текстуальных связей между источниками применяются также количественные методы анализа текста.

\textbf{\emph{Историко-правовой подход }}используется для анализа норм обычного права, судебной процедуры и способов доказательства в связи с конкретными условиями их действия. Он позволяет сопоставлять содержание нормативных памятников с зафиксированными в них и в актовом материале способами разрешения споров, принесения клятвы, свидетельства, проведения судебных испытаний и использования письменных документов.

\textbf{\emph{Историко-лингвистический анализ }}применяется к правовой лексике, речевым формулам, глаголам речи, прямой речи и иноязычным элементам в латинских и старокаталанских памятниках. Отдельно рассматриваются заимствования, цитирования и способы объяснения терминов в самом правовом тексте. \textbf{\emph{Переводоведческий}}\textbf{\emph{ анализ}} используется при сопоставлении латинского текста «Барселонских обычаев» с его старокаталанскими версиями и позволяет выявить устойчивые терминологические соответствия, а также конкретизацию, генерализацию, модуляцию, объяснительный перевод и синтаксическую переработку исходного текста.

\textbf{\emph{Историко-антропологический }}подход используется при анализе зависимости судебной процедуры от социального положения и репутации её участников. Он позволяет рассматривать клятву, свидетельские показания, ордалию и судебный поединок с учётом требований к участникам этих действий, а также исследовать значение репутации и конфессиональной принадлежности при оценке допустимости и формы доказательства.

\textbf{\emph{Количественные методы }}применяются для предварительного выявления текстуальных связей между правовыми памятниками. После сегментации и нормализации текстов возможные соответствия отбираются с использованием TF--IDF, а их близость оценивается по нескольким показателям: косинусной близости, мере Tesserae, мягкой косинусной близости и локальному выравниванию Смита---Уотермана. Результаты проходят Парето-фильтрацию и ранжирование и затем представляются в виде графов текстуальных связей. Вычислительно выявленное сходство не рассматривается как доказательство прямого заимствования: оно служит основанием для последующего историко-филологического анализа конкретных фрагментов.

Исследование опирается на\textbf{\emph{ принцип историзма}}, а также на \textbf{\emph{общеисторические методы анализа}}: историко-генетический, историко-сравнительный и историко-типологический. Их применение позволяет учитывать условия возникновения и бытования правовых памятников, прослеживать изменения норм, формул и процедур и сопоставлять разные типы источников.

\textbf{Научная новизна исследования} состоит в характеристике и уточнении характера взаимосвязей между текстами записей обычного права, выбором языковых стратегий и конкретными способами судебного доказательства внутри каталонской правовой культуры XII--XIII вв. В диссертации показано, каким текстуальные связи сводов обычного права соотносятся с языковыми особенностями (включением определённых правовых клише, переводческими решениями, заимствованиями, цитированиями и др.), а также с представлениями о клятве, свидетельстве, репутации, ордалии, судебном поединке и письменном документе. Рассмотрение этих связей позволяет точнее описать особенности бытования каталонской правовой культуры в период изменения судебной процедуры и перераспределения авторитета между обычным, королевским, каноническим правом.

\textbf{Теоретическая значимость исследования }состоит в уточнении подходов к средневековому правовому тексту как источнику по истории правовой культуры, понимаемой как система взаимосвязей между нормой, языком, процедурой, статусом, властью и формами судебного доказательства. Работа вносит вклад в изучение способов письменной фиксации обычая на нескольких взаимосвязанных уровнях: текстовом, языковом и антропологическом уровнях. На текстовом уровне рассматриваются принципы отбора и упорядочения обычая; на языковом --- терминология и формулы; на антропологическом --- способы описания статуса, репутации, процессуальной роли человека.

\textbf{Практическая значимость исследования} заключается в возможности использования его результатов в рамках учебных курсов по истории Средних веков, источниковедению, исторической антропологии, истории средневекового города, истории Пиренейского полуострова, средневековой правовой культуре, истории права и проблемам перевода правовых текстов. Выводы диссертации могут быть полезны при комментировании и переводе средневековых латинских и старокаталанских источников, поскольку работа уточняет датировки создания отдельных источников, значение ряда терминов, юридических формул и правовых клише.

\textbf{Структура диссертации }определяется проблемным принципом и соответствует цели и задачам исследования. Работа состоит из введения, раздела, посвящённого политико-правовой эволюции каталонских графств XI--XIII вв., трёх глав, заключения, библиографии и трёх приложений. В \textbf{первой главе} рассматривается вычислительная реализация метода выявления текстуальных связей между историческими корпусами и предлагается историческая интерпретация полученных данных, на основании которой реконструируется генеалогия записей обычного права в Каталонии XII--XIII вв. \textbf{Вторая глава} посвящена анализу устного, письменного и учёного регистров языка, отражённых в компиляциях обычного права, а также реконструкции, насколько это позволяют источники, процессов записи, редактирования и переработки текста средневековыми легистами. Отдельно рассматривается проблема перевода текстов обычаев с латинского языка на старокаталанский. В \textbf{третьей главе }исследуется система судебного доказательства, представленная в нормах обычного права, зафиксированных легистами, и в судебных актах. В \textbf{заключении }обобщаются результаты и формулируются основные выводы исследования. Список использованных источников и литературы приведён в разделе \textbf{Библиография}. В \textbf{приложении 1 }помещена карта Каталонии XII--XIII вв., в \textbf{приложении 2} --- словарь используемых в работе понятий и терминов, в \textbf{приложении 3} --- пример Markdown-отчёта, формируемого в ходе вычислительного анализа текстуальных связей.

\section*{Положения, выносимые на защиту}
\addcontentsline{toc}{section}{Положения, выносимые на защиту}

\textbf{1. Разработанный вычислительный метод позволяет выявлять и систематизировать текстуальные связи между отдельными фрагментами, памятниками и корпусами правовых источников. }Последовательное сочетание сегментации и нормализации текстов, нескольких мер текстуальной близости, Парето-фильтрации, ранжирования и визуализация позволяют устанавливать общую структуру распределения связей и выделять наиболее значимые соответствия для последующего источниковедческого и текстологического анализа. Историческая интерпретация выявленных соответствий проводится на основании исходных фрагментов и позволяет различать прямое цитирование, заимствование отдельных элементов, устойчивые формулы и конструкции.

\textbf{2. Текстуальные связи между исследованными правовыми памятниками распределены неравномерно и образуют устойчивую структуру}. В сопоставлении «Барселонских обычаев» с более ранними латинскими памятниками значительное место занимают связи с римской правовой традицией и «Вестготской правдой», локализованные в отдельных статьях памятника. Внутри корпуса каталонского обычного права наиболее плотные связи выявлены между «Барселонскими обычаями» и «Обычаями Тортосы». В актовом материале IX--XII вв. особенно заметны связи с «Вестготской правдой», положения которой включались в документы в форме цитат, отдельных фрагментов и переработанных формул.

\textbf{3. Письменная фиксация словесных правовых действий в каталонских источниках XII--XIII вв. предполагала их отбор, редакционную обработку и включение в установленный порядок правового действия. }Прямая речь, речевые формулы и глаголы речи закрепляли содержание обвинения, вызова и явки в суд, клятвы, распоряжения имуществом и других словесных действий, определяли последовательность процедуры и распределяли роли между её участниками. Наиболее развёрнутая система таких формул представлена в «Обычаях Тортосы». В актовом материале фиксировались формулы признания права, уступки имущества и отказа от притязаний, а также обстоятельства произнесения, восприятия и последующего письменного удостоверения слов.

\textbf{4. Сосуществование латыни и }\textbf{старокаталанского}\textbf{ языка определяло языковую организацию каталонских правовых памятников XII--XIII вв. }В латинских сводах обычного права и актовом материале местные и иноязычные элементы вводились посредством пояснительных формул, подвергались графической и морфологической адаптации либо непосредственно включались в латинский текст; отдельные обозначения связывались составителями с определённой социальной или речевой средой. В старокаталанских памятниках латинские элементы сохранялись в правовой терминологии, дефинициях и цитатах из авторитетных текстов, а также при оформлении письменного документа --- в подписях сторон и текстах печатей публичной власти.

\textbf{5. }\textbf{Старокаталанские}\textbf{ версии «Барселонских обычаев» свидетельствуют об устойчивости ряда терминологических соответствий и о вариативности переводческой переработки латинского текста. }Переводческие изменения затрагивали лексику, грамматические формы и синтаксическую организацию текста и включали конкретизацию, генерализацию, объяснительный перевод, изменение способов выражения долженствования и возможности, перестройку синтаксических конструкций и опущение отдельных элементов. Сопоставление старокаталанских версий показывает различия в выборе лексики, составе формул и объёме передаваемого материала; переводческие решения могли изменять степень конкретности и отдельные смысловые акценты исходного текста.

\textbf{6. В Каталонии XII--XIII вв. ордалии, судебный поединок, клятва, свидетельские показания, репутационные критерии и письменные документы составляли взаимосвязанную систему судебного доказательства. }Выбор и условия применения отдельных способов доказательства зависели от характера дела, процессуального положения сторон, их социального статуса и в ряде случаев конфессиональной принадлежности. На протяжении рассматриваемого периода изменялось соотношение этих способов: ограничивалась сфера применения отдельных видов ордалий, конкретизировались процессуальные функции клятвы, требования к свидетелям и допросу, значение репутационных критериев и порядок использования письменного документа.

\textbf{7. В XIII в. свидетельское показание и письменный документ включались во всё более подробно регламентированный порядок получения, удостоверения и проверки доказательств. }Оценка свидетельского показания учитывала репутацию лица, его статус, отношения со сторонами и соблюдение установленных требований к свидетелю. В «Обычаях Тортосы» репутационные критерии получили значение при определении пригодности свидетеля и проведении розыскного расследования. Доказательное значение письменного документа связывалось с публичным удостоверением, полномочиями писца или нотариуса, происхождением записи из городской канцелярии, соблюдением формуляра, наличием подписей и знаков и сохранением записи в реестре. Свидетельские показания использовались при проверке обстоятельств составления документа и восстановлении содержания утраченной записи, а розыскное расследование получило более определённые правила проведения, круга участников и фиксации полученных сведений.

\textbf{Достоверность результатов исследования }определяется комплексным характером источниковой базы, критическим анализом источников и учетом их жанровой, языковой и функциональной специфики. Обоснованность выводов обеспечивается обращением к отечественной и зарубежной историографии по истории средневекового каталонского обычного права, правового языка и судебной процедуры, а также применением взаимодополняющих историко-правовых, источниковедческих, историко-лингвистических и количественных методов.

\textbf{А}\textbf{пробация работы}. Диссертация была заслушана и рекомендована к защите на заседании кафедры истории средних веков исторического факультета МГУ имени М.В. Ломоносова N сентября 2026 г.

Основные положения диссертации были \textbf{апробированы} в докладах на научных конференциях и круглых столах: «Язык и текст: пространство применения и интерпретации» (25 апреля 2023 г., Факультет иностранных языков и регионоведения МГУ имени М. В. Ломоносова), «Власть и правосудие в Европе в Средние века и раннее Новое время: институты, идеи, практики» (10--11 октября 2023 г., Исторический факультет МГУ имени М. В. Ломоносова, ИВИ РАН), «Средневековый город: современные проблемы и их решения» (27 октября 2023 г. Библиотека имени М.Ю.~Лермонтова), «История права в Средние века и раннее Новое время: человек --- институты --- практики» (2--3 декабря 2024 г. Исторический факультет МГУ им. М.В.~Ломоносова).

Основные положения диссертационного исследования отражены в четырёх научных статьях автора, опубликованных в 2023--2026~гг. в рецензируемых научных изданиях, рекомендованных для защиты в диссертационном совете МГУ по специальности 5.6.2 «Всеобщая история».
